\subsection{Comparative evidence and diagnostics}
\label{sec:exst-comparison}

Here \term{existential} names Iridian's construction for asserting existence
or presence in a contextual domain; it does not name a universal syntactic
category. McNally's cross-linguistic survey distinguishes structural,
compositional, and discourse-functional variation \cite{mcnally2016existential}.
Sarda and Lena likewise separate the definition of existential constructions
from their semantic and discourse functions, including their polyfunctionality
\cite{sardalena2023}. These comparisons motivate keeping nominal syntax, event
entailment, and event-reporting use separate in the present account.

Tagalog provides competing analyses rather than a template for Iridian.
Sabbagh analyses an existential predicate with a nominal sole internal
argument and derives a definiteness effect from its property-denoting pivot
\cite{sabbagh2009}. This supports the availability of a nominal analysis;
it does not establish the same definiteness restriction for Iridian's
independently specified anchoring system. For a specialised Tagalog event
existential, Aldridge instead proposes a reduced \synhead{vP} complement,
using extraction as evidence against a nominal complement
\cite{aldridge2011event}. Latrouite and Van Valin contrast
\textit{May librong binili ng babae} \trsl{There is a book which the woman bought}
with \textit{May biniling libro ang babae} \trsl{The woman bought a book}:
the latter has an overt nominative actor despite undergoer-voice morphology.
They propose a reanalysis from a possessive construction containing an
internally headed relative \cite{latrouitevanvalin2014event}. These are
language-specific analyses of a morphosyntactic contrast, not consequences
of an English translation with \trsl{someone} or \trsl{something}.

For Aklanon, Wellstood proposes a relative-clause complement with optional
possessor raising and control of a description-internal actor
\cite{wellstood2022}. The proposal shows why superficially comparable
Philippine constructions need not receive a uniform clausal analysis.
Law's account of the Tagalog impersonal construction treats pluralisation,
adverbial interpretation, argument marking, locality, and extraction as
relevant evidence \cite{law2010impersonal}. Such diagnostics suggest research
questions for Iridian; they supply no missing plural paradigm or negative-only
number prohibition. Kimoto's survey of twenty northern Philippine languages
finds nominative, non-nominative, linker, and incorporation patterns for
existential or possessive themes \cite{kimoto2026existential}. Iridian's
direct-case pivot is therefore a language-specific generalisation, not a
universal result of existential semantics.

Three diagnostics would bear on any proposed additional Iridian construction:
whether event adverbials have a scope independent of the nominal description
and matrix existential domain; whether an event-internal participant acquires
a matrix grammatical relation despite the ordinary Voice and case rules; and
whether dependencies cross boundaries that otherwise block extraction from
nominals (Chapter~\ref{ch:embedding}). These are open tests of competing
analyses, not newly licensed dependencies. Until independent evidence requires
another structure, the event-reporting construal described here retains the
nominal pivot, its ordinary participant orientation, and the existing case
system. The existential--possessive relation
(\S\,\ref{sec:existential-constructions}) offers a possible comparative route
for reanalysis, not evidence that such a change has occurred in Iridian.
